
% ============================================================
% Tabel 5: Metrik Evaluasi
% ============================================================

\begin{table}[!htbp]
\centering
\caption{Metrik evaluasi komparatif \emph{GraphRAG Framework}.}
\label{tab:metrik-evaluasi}
\begin{tabularx}{\textwidth}{|l|>{\raggedright\arraybackslash}X|>{\raggedright\arraybackslash}p{3.5cm}|}
\hline
\textbf{Metrik} & \textbf{Alasan Pemilihan} & \textbf{Kondisi A/B/C} \\
\hline
\emph{Cosine Similarity} &
    Metrik utama \textcite{DaSilva2025}; memastikan komparabilitas
    \emph{apple-to-apple} &
    Semua \\
% \hline
% ROUGE-1/2/L &
%     Kemiripan berbasis \emph{n-gram} dengan jawaban \emph{ground truth} &
%     Semua \\
% \hline
% BERTScore &
%     Kemiripan semantik berbasis \emph{embedding} kontekstual &
%     Semua \\
% \hline
\emph{Hallucination Rate} &
    Persentase klaim yang tidak terdukung dalam jawaban
    (dinilai LLM-as-Judge \autocite{Gu2026}) &
    Semua \\
\hline
\emph{Attribution Coverage} &
    Persentase jawaban yang mencantumkan $\geq 1$ kutipan SO thread &
    Semua (C sebagai target 100\%) \\
\hline
\emph{LLM-as-Judge Score} &
    Penilaian kualitas holistik: \emph{faithfulness},
    \emph{comprehensiveness}, \emph{relevance} &
    Semua \\
\hline
Uji Statistik &
    UJI \emph{Wilcoxon Signed-Rank} (dua sisi, $\alpha = 0{,}05$); dipilih karena distribusi skor \textit{faithfulness} dan \textit{hallucination rate} tidak dapat diasumsikan normal pada skala terbatas $[0, 1]$ \autocite{Kamalipour2026} & A vs C, B vs C \\
\hline
\end{tabularx}
\end{table}